\documentclass[11pt]{article}
\usepackage[a4paper,margin=25mm]{geometry}
\usepackage{amsmath,amssymb,amsthm,booktabs,hyperref}
\newcommand{\E}{\mathbb E}\newcommand{\Pp}{\mathbb P}
\newcommand{\CE}{C_{\rm E}}\newcommand{\pos}[1]{(#1)_+}
\newtheorem{theorem}{Theorem}\newtheorem{lemma}[theorem]{Lemma}
\title{One fixed-target panel: an acquired strip and a confinement barrier}
\author{Bounded research note, 28 October 2026}\date{}
\begin{document}\maketitle

\section{Statement and explicit inherited premises}
Let $n\ge N=42$ and let a centered variance-one law $P$ and raw threshold
$H$ maximize the fixed-target excess
\[
 \mathfrak a=\sqrt n\{F_{n,P}(H)-\Phi(H/\sqrt n)\}-c\rho\ge0,
 \qquad c=\CE,\quad \rho=\E|X|^3.
\]
Assume $C_{n-1}\le c$ and the fixed-target contact theorem, with closed
support contact. Let $\sigma_0=\sqrt{10}$, $\gamma=\E X^3$,
$m=\E X_+$, and
\[
 J=\sum_{(u,v)}\tau_{u,v}(v-u)^2,
 \quad \tau(t)=\E[X\mathbf1_{X\ge t}],\quad
 D=(3+\sigma_0)\rho-\gamma-3J.
\]
All gaps are actual complementary intervals of the closed support.
Suppose the acquired entry has supplied
\begin{equation}\label{conf:entry}
 |X|\le96/25,\quad p=\Pp(X>0)\in[l,h]=[0.41,0.42],
 \quad\theta\le0.27,\quad nD\le42(0.7).
\end{equation}
Here $q=1-p$, the conditional means are $-a=-m/q$ and $b=m/p$,
$V=\E[X\mid\mathbf1_{X>0}]$, and $\theta=\E(X-V)^2$.
Put $\ell=a+b=m/(pq)$ and $X=\ell(Y-p)$, $Y=B_{\rm ind}+U$,
where $B_{\rm ind}=\mathbf1_{X>0}$ and $\E(U\mid B_{\rm ind})=0$.
These hypotheses refer to one law. The numerical entry table is a
conditional input here; this note does not replay its Fourier acceptance.

\begin{theorem}\label{conf:panel}
Under the stated premises,
\begin{gather*}
 \theta<0.054,\qquad d<0.069,\qquad \rho<1.218600<1.22,\\
 0.09736\le\beta:=\E[X|X|]\le0.234,
 \qquad \ell>1.9706,\qquad \lambda:=n\E U^2<0.6,\\
 \operatorname{supp}P\subset(-1.8,2.325),\qquad
 \operatorname{supp}Y\subset(-0.6,1.6),\qquad |U|<0.6.
\end{gather*}
The quantity $d=1-(A+B)/\ell$ is the normalized inward displacement,
where $(-A,B)$ is the actual central support gap. The bound on $\lambda$
uses the weighted hypothesis $nD\le42(0.7)$, not merely $D\le0.7$.
\end{theorem}

\section{A same-law moment iteration}
First $D\le0.7<\sigma_0-1$ excludes zero from the closed support:
if zero belonged to it, summing the one-sided gap bound would give
$J\le\rho$, and thus $D\ge(\sigma_0-1)\rho\ge\sigma_0-1$.
The gap endpoints satisfy $0<A\le a$, $0<B\le b$.
Set $\Delta_A=a-A$, $\Delta_B=b-B$.
On each sign class use the measure
$d\mu_\pm=(X-V)^2\mathbf1_{\{\pm X>0\}}\,dP$, with magnitude variable
$x=|X|$.

\begin{lemma}[One common resource]\label{conf:cost}
Let $p_*=(4-\sigma_0)/2$. Then
\begin{align}
 D&\ge B_0+E_0+\int k_+\,d\mu_++\int k_-\,d\mu_-,\label{conf:joint}\\
 B_0&=\frac{(6+2\sigma_0)(1-\theta)^{3/2}}{\sqrt{pq}}(p-p_*)^2,
 \quad E_0=6m\{(A+B)(\Delta_A+\Delta_B)+\Delta_A\Delta_B\},\notag\\
 k_+(x)&=(\sigma_0-1)(x+2b)+6B,
 \quad k_-(x)=(\sigma_0+1)(x+2a)+6A.\notag
\end{align}
All four terms on the right are nonnegative. Moreover,
\begin{equation}\label{conf:product}
 AB\ge\frac{\sigma_0-1-D}{3},\qquad
 E_0\ge6m\ell^2(1-d)d.
\end{equation}
\end{lemma}
\begin{proof}
For positive outer gaps, comparison with
$2\int(t-B)\tau(t)\,dt$ gives
$J_o^+\le T_+-2BR_+$, where
$T_+=\E X_+^3-mB^2$ and $R_+=\E X_+^2-mB\ge0$.
Reflection gives $J_o^-\le T_--2AR_-$, with the analogous negative
quantities. Add the central contribution $m(A+B)^2$ and substitute in D:
\[
 D\ge(\sigma_0-1)\E X_+^3+(\sigma_0+1)\E X_-^3
            -6mAB+6BR_++6AR_-.
\]
Conditional centering gives
\[
 \E X_+^3=mb^2+\int(x+2b)\,d\mu_+,
 \quad R_+=\mu_+(\mathbb R)+m\Delta_B,
\]
and the corresponding negative identities. The constant part is $B_0+E_0$;
use $a=\ell p$, $b=\ell q$, and
\[
 (\sigma_0+1)p^2+(\sigma_0-1)q^2-6pq
                  =(6+2\sigma_0)(p-p_*)^2.
\]
This proves \eqref{conf:joint}. The simpler gap bound
$J\le\rho+2mAB$ gives
$6mAB\ge\sigma_0\rho-\gamma-D\ge\sigma_0-1-D$.
Since $m\le1/2$, the first part of \eqref{conf:product} follows.
The second uses $A+B=\ell(1-d)$ and
$\Delta_A+\Delta_B=\ell d$.
\end{proof}

Here is the precise finite iteration, with $D_*=0.7$ and
$s_*=3.162277660<\sigma_0$. Write
$v_-=l(1-l)$ and $v_+=h(1-h)$.
Initially
\[
 a\le a^+=\sqrt{h/(1-h)},\quad b\le b^+=\sqrt{(1-l)/l},\quad
 P_*=(s_*-1-D_*)/3>0.
\]
Take outward bounds
$A_-=P_*/b^+$, $B_-=P_*/a^+$ and
$c_-=2\sqrt{P_*v_-}\le(A+B)/\ell$, with $\Theta=0.27$.
Given these bounds, let
\[
 a_-^m=\sqrt{(1-\Theta)l/(1-l)},\quad
 b_-^m=\sqrt{(1-\Theta)(1-h)/h},
\]
and set
\begin{align*}
 K_\theta&=\min\{(s_*-1)(B_-+2b_-^m)+6B_-,
                    (s_*+1)(A_-+2a_-^m)+6A_-\},\\
 \Theta'&=\min\{\Theta,D_*/K_\theta\},\qquad
 e_- =\sqrt{(1-\Theta')/v_+},\quad m_- =\sqrt{(1-\Theta')v_-},\\
 d'&=\min\{1-c_-,D_*/(6m_-e_-^2c_-)\},\quad c_-'=\max\{c_-,1-d'\},\\
 A_-'&=\max\{A_-,e_-(l-d')\},\quad
 B_-'=\max\{B_-,e_-(1-h-d')\}.
\end{align*}
Every square root in the checker is rounded in the required direction.
The first update follows from $D\ge K_\theta\theta$ in
\eqref{conf:joint}; the second follows from \eqref{conf:product}.
Only already proved bounds are used before each update. Four repetitions give
the following strict upper bounds:
\begin{center}
\begin{tabular}{ccc}\toprule
step&$\theta$&$d$\\\midrule
1&0.077629&0.094983\\
2&0.055626&0.069572\\
3&0.053266&0.067419\\
4&0.053060&0.067242\\\bottomrule
\end{tabular}
\end{center}
Thus the simpler closed bounds $\Theta=0.054$, $d_0=0.069$ are valid.
They imply $\ell>E=1.9706$, and hence
$A\ge E(l-d_0)=0.6719746$, $B\ge E(1-h-d_0)=1.0069766$.

To retain coupling in the third moment, let
$\rho_B(p)=(1-2pq)/\sqrt{pq}$ and
$\delta\rho=\rho-(1-\theta)^{3/2}\rho_B(p)$.
The same conditional expansion gives
$\delta\rho=\sum_\pm\int(x+2a_\pm)\,d\mu_\pm$.
With $R=96/25$ and $L=R+2b^+$, put
\[
 k=\min\left\{s_*-1+\frac{6B_-}{R+2b^+},
                 s_*+1+\frac{6A_-}{R+2a^+}\right\}>0.
\]
Since $p<1/2$, $b^+\ge a^+$, and
\[
                 \delta\rho\le\min\{L\theta,D_*/k\}.
\]
Take an outward upper $r_B\ge\rho_B(l)$ and $a_*=29/20$.
The gate $a_*^2<(9/4)(1-\Theta)$ gives
$(1-\theta)^{3/2}\le1-a_*\theta$ for $0\le\theta\le\Theta$.
Also $L-a_*r_B>0$. The piecewise linear maximum therefore gives
\[
 \rho\le r_B+(L-a_*r_B)\min\{\Theta,D_*/(kL)\}
                 <1.218600<1.22.
\]
The moment bound uses the original radius $R$; it is not based on the
support improvement proved next. The signed second moment identity gives
\[
 1-2h-2(1-h)\Theta\le\beta\le1-2l+\Theta,
\]
which is the claimed interval. Finally the refreshed cost satisfies
$K_\theta>13.16206$. Multiplying $D\ge K_\theta\theta$ by $n$ and using
$\lambda=npq\theta/(1-\theta)$ gives
\[
 \lambda\le\frac{42D_*v_+}{(1-\Theta)K_\theta}<0.575188<0.6.
\]

\section{Signed contact rays acquire the initial strip}
Fixed-target contact and $C_{n-1}\le c$ imply at every support point $x$
\begin{equation}\label{conf:contact}
 0\le\mathcal H_n(H,x)+c\{-|x|^3+\tfrac32\rho x^2+3\beta x+\rho e_n\}
                         -n\mathfrak a,
\end{equation}
where $e_n=n(\sqrt{n/(n-1)}-1)-1/2<8/875$ and $\mathcal H_n$ is the
Gaussian heat remainder. The final term is nonpositive and can be dropped.

The retained fourth-order heat expansion remains valid for $u>\sqrt3$,
even though its earlier simple envelope was stated only for $u\ge5/2$.
With $v=n-1$, $s=\sqrt{n/v}$, $z=H/\sqrt v$,
\[
 \mathcal H_n(H,\pm u)/s^3
 =\pm A_u(1-z^2)\phi(z)+B_u\phi'''(z)+R_u,
\]
where $A_u=(u^3-3u)/6>0$,
$B_u=(u^4-6u^2+3)/(24\sqrt v)$, and
\[
 |R_u|\le\frac{u^5/100+3u/20}{v}
                +\frac{51(3u^2-2)}{960v^{3/2}}.
\]
The Taylor remainder derivation uses $u^2-1\ge0$; this holds on the
present enlarged domain. Only the sign of $B_u$ changes, so use
$|B_u|$ in the bound $\|\phi'''\|_\infty<2\phi(0)$.
With $S=81/80$, $c_-=0.4095<c$, define
\begin{align*}
 E_u&=u^5/4100+3u/820+17(3u^2-2)/78720,\\
 Q_{r,b}(u)&=-u^3+\tfrac32ru^2+3bu+(8/875)r,\\
 P_{r,b}(u)&=\tfrac25S^3\{A_u+|u^4-6u^2+3|/72\}
                                  +S^3E_u+c_-Q_{r,b}(u).
\end{align*}
The sign $Q<0$ is checked before using the lower multiplier $c_-$.
For the negative ray choose $b=-\beta_-$, and for the positive ray
choose $b=\beta_+$. With $r=1.22$, exact rational arithmetic gives
\[
 P_{r,-0.09736}(1.8)<-0.093411,
 \qquad P_{r,0.234}(2.325)<-0.005327.
\]
To cover the entire rays, define
$L=(2/5)(u^2/2+1/2+8/875)+c_-Q'(u)$.
The uniform derivative estimate
$|\partial_x\mathcal H_n|\le(2/5)(x^2/2+1/2+e_n)$
and the six strict signs $Q,Q',Q'',L,L',L''<0$ at each boundary suffice:
$Q'''=-6$ and $L''=2/5-6c_-<0$ propagate the signs.
The checker verifies all twelve signs. Hence \eqref{conf:contact}
excludes $x\le-1.8$ and $x\ge2.325$, including boundary atoms.

The exact inequalities
\[
 1.8<E(l+0.6),\qquad 2.325<E(1+0.6-h)
\]
now give $-0.6<Y<1.6$. The inner endpoints give $U_0\le d<0.069$
and $U_1\ge-d>-0.069$; the outer strip gives the other two sides.
Thus $|U|<0.6$. This completes Theorem~\ref{conf:panel}, including zero
noise. In that case both conditional variances vanish and $U=0$ exactly;
no division by noise or limiting argument is used.

\section{What transfers to the next confinement step}
In Bernoulli coordinates the fixed-target contact majorant obeys
\[
 \sqrt n(-Q_c'(y))=\ell\phi(H/\sqrt n)
 +\frac\ell n\{3c\beta+(H\phi(H/\sqrt n)+3c\rho)t-3ct|t|\},
 \quad t=\ell(y-p).
\]
This is exactly the retained derivative with coefficient $c$.
The extra excess affects an additive constant and disappears under
differentiation. It does not license replacing the actual deleted sums
by score bounds without proving those separate bounds.

The present attempt instead uses strict-gap antichain concentration.
It needs no deleted-score constants, but it does retain both actual
indices $n-1,n-2$, the latter at least $40$. Put
\[
 r_L=l-(1-l)/40,\quad r_R=h+h/40,\quad
 v_* =\min\{r_L(1-r_L),r_R(1-r_R)\}.
\]
Finite Stirling gives a valid cap $M$ whenever
\[
 M^2>\frac{0.16}{v_*(1-1/480)^2},\qquad
 F_d^2>\max\{(42/41)^3,42/40\}.
\]
The checker uses a directed $M$ and $F_d=1.037$.
The coefficient $c_\rho=24$ is valid here because
$24^2[l(1-l)]^3>[3(1-l)+0.6]^2$; no old strip width is substituted.

For a target amplitude $r_0=1/2$, choose both buffers $\tau=1/5$.
The necessary geometry is strict:
\[
 1-0.6-2d_0>0,\quad d_0+2\tau<r_0,\quad
 1-d_0-\tau>0,\quad1-2d_0-2\tau>0.
\]
The lower interval widths also remain strictly below $1-d_0$ over
the whole acquired strip. The old particular choice $7r_0/15$ for
the upper buffer fails $d_0+2\tau<r_0$ and is not reused.

For this diagnostic the frozen entry row additionally supplies
$\mathcal R_c< T=2.016$ and $|H|<1.814$.
The raw license $T<42c_-$ gives the single conserved resource
\[
 (n+H^2)\mathfrak a+n\phi(0)D/6+n^{3/2}L_{\rm leak}
 <\mathcal B(H):=T-\eta H^2,
 \qquad \eta=(c_- -T/42)/2>0.
\]
This follows from the fixed-target Gaussian phase-price lemma; a defect
maximum alone would not give this joint resource. Every term refers to
the same law and raw threshold. The strict inequality follows from the
outward raw ceiling, not from altering an acceptance target.

The bounded check evaluates the retained slope and resource formulas on
eight closed phase intervals, eight closed noise intervals in $[0,0.6]$,
and both signs of $H$: 128 boxes. Four are excluded by the resource; 195
of the remaining 248 ray gates fail. Thus this invocation does not prove
$|U|<1/2$. Its zero-noise case is independently trivial; the failed
positive-noise gates cannot be removed on that account.

\section{A buffer-independent obstruction to this sufficient bound}
The failure is not merely caused by the coarse subdivision.
Fix the singleton $H=5/4$, $\lambda=3/20$ in the same panel and retain
the preceding outward moment, scale and concentration constants.
The upper contact side is
$L_g-0.6U_+<0.216059$.
Let $b_*$ be the retained positive coefficient of $\theta$ in the
common moment resource, and put
\[
 R_* =\mathcal B(H)/\kappa-b_*E^2\lambda>0,
 \qquad\kappa=997/15000.
\]
The upper geometry permits only
$0<\tau<(r_0-d_0)/2=0.2155=:\tau_*$.
The directed-distance term is affine in this entire interval:
\begin{align*}
 A(\tau)&=A_0+A_1\tau,\\
 A_0&=F_dM\left\{\frac{d_0+H/E-l}{l}
                   +\frac{1+d_0+H/E-l}{1-h}\right\},\\
 A_1&=F_dM\{1/l+2/(1-h)\}.
\end{align*}
Both distance numerators are positive even at $\tau=0$.
The old resource error majorant has coefficient
$z\ge\kappa/((1-\Theta)\tau)$.
Every tangent majorant also contains the leakage vertex
$v_J=\kappa R_*/(42(1-\Theta)\tau)$, whose displacement is zero;
its remaining tangent intercept is nonnegative.
Consequently \emph{each of those error majorants}, for every allowed
buffer and every nonnegative tangent parameter, is at least
\[
 f(\tau)=\frac{A_0}{42}+\frac{A_1\tau}{42}+\frac K\tau,
 \qquad K=\frac{\kappa R_*}{42(1-\Theta)}.
\]
This lower-bounds the error bounds, not the true probabilistic error.
The exact gate $A_1\tau_*^2/42<K$ makes $f$ strictly decreasing on
$(0,\tau_*]$, so every admissible choice exceeds
\[
 f(\tau_*)>0.23161994,
 \qquad f(\tau_*)-(L_g-0.6U_+)>0.01556155.
\]
Thus no refinement of the phase/noise mesh at this singleton, no choice
of its upper buffer, and no additional tangent parameter repairs this
particular sufficient-bound attempt with these fixed envelopes.
It is not a counterexample to confinement or to the sharp conjecture.

The current phase resource is $\mathcal B(5/4)=1.733578125$.
Keeping the other envelopes, a necessary condition for removing this
leakage-vertex obstruction is a resource ceiling below
\[
 \kappa b_*E^2\lambda+
 42(1-\Theta)\tau_*\{L_g-0.6U_+-A(\tau_*)/42\}
 <1.600336284.
\]
The exact rational threshold is recorded by the checker. Passing this
threshold alone would still not prove the remaining confinement gates.
A coherent phase quadratic for the same active moment cell is a possible
next input; separate favorable maxima of $D,H$ do not supply it.

\section{Evidence and scope}
The standard-library exact checker binds one frozen probability-panel
object, verifies four moment updates, two signed boundary inequalities,
twelve whole-ray signs, the noise bound, all listed geometry gates, the
128 diagnostic boxes, and the continuous buffer obstruction. Eight
negative controls reject missing/reordered source panels, a changed
defect input, duplicate keys, a false smaller support radius, the wrong
negative-ray beta sign, the invalid old buffer, and a false complete
half-amplitude claim. No inherited numerical program is executed.

Mathematical status: Theorem~\ref{conf:panel} is conditional on the
fixed-target contact, sharp predecessor and acquired entry premises.
The next half-amplitude invocation remains unproved, with an exact
obstruction to the displayed sufficient bounds. Execution status:
fresh rational checks of the new component and diagnostic; no native
INITIAL42/LOCAL42 acceptance, finite base or global theorem is claimed.
\end{document}
